% Propietario conservado: /Users/ruben/Documents/New project/output/REVISION_ARGUMENTAL_APERTURAS_CIERRES_20260915/VII/delivery/MANUSCRIPT_VII_ES_EN_COMPLETE/source_es/manuscrito/32_eliminacion_torsional_no_lineal.tex
% FORMALIZACION_NUEVA / CERTIFICADO_NUEVO, 10-09-2026.
% Composición de las operaciones efectivas ya reunidas en 30e y 31.
% No modifica la selección de rutas, la frontera, los acoplamientos ni
% la normalización de la acción; no evalúa una amplitud cosmológica s_0.
% Propietarios y alcance: desarrollo/COMPOSICION_METRICA_REV06.md.
% Control independiente: pruebas/verificar_eliminacion_torsional_no_lineal.py.

\section{Élimination torsionnelle et réponse métrique de l'extension minimale}
\label{vii:vii:sec:eliminacion-torsional-no-lineal}

L'extension minimale fournit un champ, un poids de réponse et un courant au moyen d'une seule opération variationnelle. Son incidence transportée dépend de la connexion à travers les produits de route de \eqref{vii:vii:eq:derivada-incidencia-rutas}. Cette dépendance est conservée lors de la composition du courant avec l'équation de Cartan--Holst. L'élimination de la connexion s'effectue donc sur l'action complète de cette extension. Le cas affine de \ref{vii:vii:thm:respuesta-cuadratica} est inclus comme une spécialisation exacte de la construction suivante.

\subsection{Composition variationnelle à un niveau cellulaire}
\label{vii:vii:subsec:composicion-torsional-celular}

Fixons un niveau fini, ses routes, la représentation de transport et
l'appariement cellulaire de
\ref{vii:vii:thm:componentes-minimo-conexion}. Soient \(p\) des coordonnées
réelles de la cotétrade et des champs matériels dans une trivialisation
locale déclarée, et \(k\in\mathbb R^d\) les coordonnées de la
contorsion \(K\). La famille \(p\mapsto e(p)\) conserve la
non-dégénérescence de la cotétrade. La connexion qui agit dans chaque route est
\begin{equation}
 \omega(p,k)=\mathring\omega(e(p))+K(k),\qquad
 F(p,k)=S_{\min}[e(p),\omega(p,k),\Psi(p)]
       =\mathfrak a(p,k)E(p,k).
 \label{vii:vii:eq:accion-minima-en-contorsion}
\end{equation}
Les espaces et les bases de cette expression sont transportés vers des espaces
fixes avant de dériver. On travaille dans un ouvert où l'incidence
\(\mathsf C(p,k)\) est surjective et les familles sont \(C^2\).
Cette régularité supplémentaire permet d'examiner le hessien de l'action.
Les couplages \(G_{\rm int},c_{\rm int},\gamma_{\rm HMT}\)
conservent la carte et la convention de variation de
\ref{vii:vii:subsec:dominio-corriente-cartan}.

Soit \(Q(p,k)\) l'intégrale, ou la réalisation cellulaire déclarée, de
la forme quadratique
\eqref{vii:vii:eq:forma-cuadratica-contorsion}. Dans des bases réelles,
\begin{equation}
 Q(p,k)=\tfrac12 k^{\mathsf T}\mathsf A(p)k,
 \qquad \mathsf A(p)^{\mathsf T}=\mathsf A(p).
 \label{vii:vii:eq:cuadratica-celular-contorsion}
\end{equation}
L'opérateur \(\mathsf A\) provient de l'action gravitationnelle et de son accouplement, tandis que le gramien \(\mathsf C\mathsf C^*\) provient du problème d'extension. La positivité du second garantit l'existence du champ minimisant ; la forme gravitationnelle conserve sa propre signature. Les termes de bord sont maintenus selon \eqref{vii:vii:eq:borde-contorsion} ; les formules intérieures suivantes s'appliquent avec des variations admissibles qui annulent leur contribution, ou avec la complétion de bord correspondante.

Définissons le covecteur de réponse par
\begin{equation}
 d_kF(p,k)[h]= -\tfrac12 h^{\mathsf T}s(p,k),
 \qquad s(p,k)=\mathsf M_{\rm cel}(p,k)\sigma(p,k).
 \label{vii:vii:eq:covector-torsional-minimo}
\end{equation}
Ici, \(s\) et \(\sigma\) conservent les types distincts de
\eqref{vii:vii:eq:componentes-minimo-conexion} : le premier contient les
coefficients de la différentielle et le second les coordonnées du
courant dans l'appariement matériel. Si l'on change de base dans
\(k\), cet appariement est également transporté. Lorsqu'on différencie
\(s=\mathsf M_{\rm cel}\sigma\), on inclut les dérivées des
deux facteurs lorsque l'appariement dépend du paramètre.

\begin{proposition}[Équation composée de l'extension et de la contorsion]
\label{vii:vii:prop:ecuacion-torsional-no-lineal}
La stationnarité intérieure de l'action par rapport à la contorsion est
\begin{equation}
 \mathcal R(p,k):=\mathsf A(p)k-\tfrac12s(p,k)=0.
 \label{vii:vii:eq:ecuacion-torsional-no-lineal}
\end{equation}
Chaque composante de \(s\) est calculée à partir de la même extension :
\begin{equation}
 \begin{split}
 s_j={}&4\mathfrak a\chi\operatorname{Re}
       \langle y,(\partial_{k_j}\mathsf C)f_{\min}
                         -\partial_{k_j}b\rangle\\
 &-2\mathfrak a(\partial_{k_j}\chi)\langle b,y\rangle
       -2(\partial_{k_j}\mathfrak a)E.
 \end{split}
 \label{vii:vii:eq:seleccion-variacional-corriente-k}
\end{equation}
Lorsque l'appariement cellulaire réalise l'équation de connexion de
\ref{vii:vii:thm:ecuacion-cartan-holst}, cette équation équivaut à composer
\(\sigma(p,k)\) avec l'inverse de Holst, l'inverse de
\(\mathsf L_e\) et la reconstruction de la contorsion :
\begin{equation}
 K=\mathsf C_e^{\rm tor}\,\mathsf L_e^{-1}
           \bigl(\kappa c\,\mathsf P_\gamma^{-1}
             \sigma[e,\mathring\omega(e)+K,\Psi]\bigr).
 \label{vii:vii:eq:composicion-implicita-cartan}
\end{equation}
Le symbole \(\mathsf C_e^{\rm tor}\) désigne exclusivement l'application \(T\mapsto K\) de \eqref{vii:vii:eq:contorsion-componentes}, à cotétrade fixe.
\end{proposition}

\begin{proof}
La différentielle de \(Q+F\), dans la direction arbitraire \(h\), est
\[
 d_k(Q+F)[h]
   =h^{\mathsf T}\bigl(\mathsf A k-\tfrac12s\bigr).
\]
Son annulation pour toute direction prouve \eqref{vii:vii:eq:ecuacion-torsional-no-lineal}. Le théorème~\ref{vii:vii:thm:componentes-minimo-conexion}, appliqué à \(\partial_{k_j}\), donne \eqref{vii:vii:eq:seleccion-variacional-corriente-k}. Avec \(p\) fixe, \(\delta\omega=\delta K\). L'équation de Cartan se transforme successivement au moyen de \eqref{vii:vii:eq:inversa-holst}, \eqref{vii:vii:eq:torsion-explicita} et \eqref{vii:vii:eq:contorsion-componentes} ; chacune de ces opérations est inversible dans son domaine. Cela prouve l'équivalence avec \eqref{vii:vii:eq:composicion-implicita-cartan} lorsque les deux expressions utilisent la même réalisation matérielle et cellulaire.
\end{proof}

La composition conserve la connexion dans les deux membres. Les routes
qui parcourent plusieurs cellules peuvent coupler les composantes du
courant entre elles. La reconstruction ponctuelle de \(T\) à partir
d'un courant donné et la résolution conjointe de
\eqref{vii:vii:eq:ecuacion-torsional-no-lineal} sont, dans ce cas,
des opérations successives. Cette formulation finie maintient explicite
la représentation des routes qui intervient dans la fonctionnelle.

\subsection{Existence locale d'une branche stationnaire}
\label{vii:vii:subsec:rama-torsional-local}

\begin{theorem}[Continuation locale de la solution torsionnelle]
\label{vii:vii:thm:rama-torsional-local}
Supposons que \((p_0,k_0)\) appartient à l'ouvert précédent,
\(\mathcal R(p_0,k_0)=0\), et que le hessien total
\begin{equation}
 \mathsf H_0
 =\mathsf A(p_0)+\partial_k^2F(p_0,k_0)
 =\mathsf A(p_0)-\tfrac12\partial_ks(p_0,k_0)
 \label{vii:vii:eq:hessiano-total-torsional}
\end{equation}
est inversible. Il existe des voisinages de \(p_0\) et de \(k_0\) dans lesquels
une unique branche \(C^1\), \(k=k_*(p)\), satisfait
\eqref{vii:vii:eq:ecuacion-torsional-no-lineal} et passe par \(k_0\).
Sur cette branche,
\begin{equation}
 \partial_{p_i}k_*
 =-\bigl[\mathsf A+\partial_k^2F\bigr]^{-1}
     \left[(\partial_{p_i}\mathsf A)k_*
                      +\partial_{p_i}\partial_kF\right].
 \label{vii:vii:eq:derivada-rama-torsional}
\end{equation}
\end{theorem}

\begin{proof}
L'inversion du gramien est \(C^2\) sur l'ouvert de rang maximal. Ainsi, \(F\) est \(C^2\) et \(\mathcal R\) est \(C^1\). Sa dérivée par rapport à \(k\), en \((p_0,k_0)\), est \(\mathsf H_0\). Le théorème des fonctions implicites fournit la branche, sa régularité et son unicité au voisinage du point. On peut également voir l'existence au moyen de l'application
\[
 \mathcal T_p(k)=k-\mathsf H_0^{-1}\mathcal R(p,k).
\]
Sa dérivée par rapport à \(k\) s'annule en \((p_0,k_0)\).
Dans une boule suffisamment petite, sa norme est au plus
\(q<1\). En réduisant le voisinage de \(p_0\), on obtient
\(\|\mathcal T_p(k_0)-k_0\|<(1-q)r\), où \(r\) est le
rayon de la boule. Ainsi, \(\mathcal T_p\) envoie la boule fermée
dans elle-même et est contractante. Son point fixe est la solution locale.
Différentier \(\mathcal R(p,k_*(p))=0\) et isoler
\(\partial_{p_i}k_*\) donne
\eqref{vii:vii:eq:derivada-rama-torsional}.
\end{proof}

L'hypothèse se vérifie sur l'action composée : le rang du
gramien assure le minimum d'écran et le hessien total contrôle
la branche torsionnelle. L'unicité affirmée est locale autour de la
solution spécifiée ; d'autres branches peuvent coexister dans le même
domaine de rang plein.

\subsection{Action effective et contribution non affine}
\label{vii:vii:subsec:accion-efectiva-no-afin}

\begin{theorem}[Identité exacte d'élimination]
\label{vii:vii:thm:eliminacion-torsional-no-lineal}
Soit \(k_*(p)\) une branche stationnaire. Supposons en outre que le
segment \(t\mapsto (p,tk_*)\), \(0\leq t\leq1\), demeure
dans l'ouvert où \(F\) est défini. La correction d'action
par rapport à la réalisation de Levi--Civita est
\begin{equation}
 \begin{split}
 \Delta S(p)
  &:=Q(p,k_*)+F(p,k_*)-F(p,0)\\
  &=\tfrac14 k_*^{\mathsf T}s(p,k_*)
     -\tfrac12\int_0^1
           k_*^{\mathsf T}s(p,tk_*)\,dt.
 \end{split}
 \label{vii:vii:eq:accion-efectiva-no-lineal}
\end{equation}
De manière équivalente,
\begin{equation}
 \Delta S=-\tfrac14 k_*^{\mathsf T}s(p,k_*)
  +\tfrac12\int_0^1 k_*^{\mathsf T}
       \bigl[s(p,k_*)-s(p,tk_*)\bigr]\,dt.
 \label{vii:vii:eq:correccion-no-afin-integrada}
\end{equation}
Ces identités se rapportent à l'action intégrée ou cellulaire. Dans une réalisation à densité locale, elles s'appliquent à la densité avec son accouplement de formes correspondant. L'action effective conserve séparément le terme de bord déclaré.
\end{theorem}

\begin{proof}
Par la règle de la chaîne et
\eqref{vii:vii:eq:covector-torsional-minimo},
\[
 \frac{d}{dt}F(p,tk_*)
      =-\tfrac12 k_*^{\mathsf T}s(p,tk_*).
\]
Intégrer entre \(0\) et \(1\) donne la différence des actions
matérielles. L'équation de stationnarité, évaluée dans la
direction \(k_*\), et l'homogénéité quadratique de \(Q\) donnent
\[
 2Q(p,k_*)=\tfrac12 k_*^{\mathsf T}s(p,k_*).
\]
La somme des deux identités démontre
\eqref{vii:vii:eq:accion-efectiva-no-lineal} ; ajouter et soustraire
\(\tfrac12k_*^{\mathsf T}s(p,k_*)\) fournit
\eqref{vii:vii:eq:correccion-no-afin-integrada}.
\end{proof}

\begin{corollary}[Spécialisation affine]
\label{vii:vii:cor:recuperacion-eliminacion-afin}
Si \(s(p,k)=s(p)\), alors
\[
 \Delta S=-\tfrac14 k_*^{\mathsf T}s(p)=-Q(p,k_*).
\]
C'est la réalisation intégrée de \eqref{vii:vii:eq:accion-efectiva-spin}.
\end{corollary}
\begin{proof}
La différence dans l'intégrale de
\eqref{vii:vii:eq:correccion-no-afin-integrada} s'annule. La seconde
égalité procède de l'identité radiale de stationnarité.
\end{proof}

\subsection{Différentielle métrique de l'action après élimination}
\label{vii:vii:subsec:diferencial-metrico-no-lineal}

\begin{theorem}[Réponse des variables géométriques]
\label{vii:vii:thm:envolvente-metrica}
Sur la branche du théorème~\ref{vii:vii:thm:rama-torsional-local},
\begin{equation}
 \partial_{p_i}\Delta S
 =\tfrac12 k_*^{\mathsf T}(\partial_{p_i}\mathsf A)k_*
  +(\partial_{p_i}F)(p,k_*)-(\partial_{p_i}F)(p,0),
 \label{vii:vii:eq:envolvente-metrica-no-lineal}
\end{equation}
où les dérivées partielles maintiennent fixes les coordonnées \(k\).
Dans des produits hermitiens fixes, chaque dérivée matérielle se calcule par
\begin{equation}
 \begin{split}
 \partial_{p_i}F={}&(\partial_{p_i}\mathfrak a)E
       +\mathfrak a(\partial_{p_i}\chi)\langle b,y\rangle\\
   &+2\mathfrak a\chi\operatorname{Re}
       \langle y,\partial_{p_i}b
                    -(\partial_{p_i}\mathsf C)f_{\min}\rangle.
 \end{split}
 \label{vii:vii:eq:diferencial-geometrico-minimo}
\end{equation}
Si \(p_i\) fait varier la cotétrade, \(\partial_{p_i}\mathsf C\) inclut la variation de \(\mathring\omega(e(p))+K(k)\), des routes réalisées et des trivialisations utilisées, selon leurs dépendances déclarées.
\end{theorem}

\begin{proof}
La dérivée totale de \(Q(p,k_*)+F(p,k_*)\) contient, outre
ses dérivées partielles, le terme
\[
 \bigl[\partial_kQ(p,k_*)+\partial_kF(p,k_*)\bigr]
                      \partial_{p_i}k_*.
\]
Le crochet est nul par l'équation de connexion. Soustraire la dérivée
de \(F(p,0)\) prouve
\eqref{vii:vii:eq:envolvente-metrica-no-lineal}.
La règle du produit et
\eqref{vii:vii:eq:diferencial-minimo-completo} prouvent
\eqref{vii:vii:eq:diferencial-geometrico-minimo}. La composition
\eqref{vii:vii:eq:accion-minima-en-contorsion} impose de différencier
la connexion de Levi--Civita lorsqu'on fait varier \(e\) à \(k\) fixé.
\end{proof}

Lorsque le produit positif d'extension varie également, sa
contribution est conservée explicitement. En coordonnées, soit
\(\mathsf h(p,k)>0\) sa matrice hermitienne et considérons le minimum
de \(f^*\mathsf h f\) sous \(\mathsf C f=b\). Alors
\begin{equation}
 \mathsf L=\mathsf C\mathsf h^{-1}\mathsf C^*,\quad
 y=\mathsf L^{-1}b,\quad
 f_{\min}=\mathsf h^{-1}\mathsf C^*y,\quad
 E=\chi b^*y.
 \label{vii:vii:eq:minimo-producto-variable}
\end{equation}
Ici, \(y\) représente le covecteur de la contrainte dans des bases fixes. La matrice positive \(\mathsf h\) conserve son rôle énergétique, distinct de celui de la métrique lorentzienne.

\begin{lemma}[Variation du produit d'extension]
\label{vii:vii:lem:producto-variable-extension}
Avec les conventions précédentes,
\begin{equation}
 \delta E=(\delta\chi)b^*y
   +2\chi\operatorname{Re}\bigl[y^*(\delta b-\delta\mathsf C f_{\min})\bigr]
   +\chi f_{\min}^*(\delta\mathsf h)f_{\min}.
 \label{vii:vii:eq:diferencial-producto-variable}
\end{equation}
\end{lemma}
\begin{proof}
La formule \(\delta\mathsf h^{-1}
=-\mathsf h^{-1}(\delta\mathsf h)\mathsf h^{-1}\) donne
\[
 \delta\mathsf L=(\delta\mathsf C)\mathsf h^{-1}\mathsf C^*
 +\mathsf C\mathsf h^{-1}(\delta\mathsf C)^*
 -\mathsf C\mathsf h^{-1}(\delta\mathsf h)
                              \mathsf h^{-1}\mathsf C^*.
\]
Lors de la différentiation de \(E=\chi b^*\mathsf L^{-1}b\), le terme \(-\chi y^*(\delta\mathsf L)y\) produit les deux termes avec \(\delta\mathsf C\) et le terme positif avec \(\delta\mathsf h\). Cela prouve l'identité complète.
\end{proof}

Dans une réalisation matérielle covariante dont la variation admet l'
appariement métrique de
\eqref{vii:vii:eq:fuentes-normalizadas-accion}, la correction détermine
\begin{equation}
 \delta_g\Delta S
 =-\tfrac12\int_M\operatorname{vol}_e\,
          \mathcal U^{S,\min}_{\mu\nu}\,\delta g^{\mu\nu},
 \qquad U^{\rm phys,\min}_{\mu\nu}
                      =c\mathcal U^{S,\min}_{\mu\nu}.
 \label{vii:vii:eq:fuente-metrica-minimo-no-lineal}
\end{equation}
Les équations \eqref{vii:vii:eq:envolvente-metrica-no-lineal} et
\eqref{vii:vii:eq:diferencial-geometrico-minimo}, complétées par
\eqref{vii:vii:eq:diferencial-producto-variable} le cas échéant,
évaluent leurs coefficients dans les variations métriques déclarées.
Le minimum d'écran, la solution torsionnelle et la dérivée métrique
sont trois opérations consécutives sur la même action.
L'action réduite conserve les éventuels couplages entre
cellules présents dans les routes originales.

La spécialisation homogène de
\ref{vii:vii:sec:dilucion-rebote} utilise en outre la réalisation du
tenseur obtenu : sa symétrie spatiale, le signe de la contraction
avec l'observateur et la loi d'échelle. Lorsque cette évaluation donne
\(\rho_{\rm corr}(a)=-s_0a^{-6}\), l'amplitude se lit dans l'
état de référence sous la forme
\(s_0=-a_{\rm ref}^{6}\rho_{\rm corr}(a_{\rm ref})\).
Cette lecture conserve l'état matériel et sa normalisation ; le
théorème d'élimination fournit l'action à contracter.

\subsection{Modèle exact de vérification non affine}
\label{vii:vii:subsec:control-no-afin}

Le modèle fini suivant vérifie les coefficients et la dépendance non affine des identités. Ses paramètres appartiennent à l'exemple algébrique et maintiennent séparé le problème de réalisation matérielle des routes.

\begin{proposition}[Élimination exacte d'une extension de rang un]
\label{vii:vii:prop:ejemplo-no-afin}
Soient \(H=\mathbb R^2\), \(Q=\mathbb R\), \(\mathsf C(k)=(1\ \ k)\), \(b=1\), \(\chi=1\), \(\mathfrak a=\lambda>0\) et \(Q_{\rm ex}(k)=k^2/2\). Alors
\begin{equation}
 f_{\min}=\frac{(1,k)^{\mathsf T}}{1+k^2},\quad
 F(\lambda,k)=\frac{\lambda}{1+k^2},\quad
 s(\lambda,k)=\frac{4\lambda k}{(1+k^2)^2},\quad
 \mathcal R(\lambda,k)
    =k\left(1-\frac{2\lambda}{(1+k^2)^2}\right).
 \label{vii:vii:eq:ejemplo-minimo-no-afin}
\end{equation}
Pour \(\lambda=25/2\), les solutions stationnaires sont
\(0,2,-2\). En \(k_*=2\),
\begin{equation}
 \mathsf H_*=\frac{16}{5},\quad s_*=4,\quad
 Q_{\rm ex}(k_*)=2,\quad F(\lambda,k_*)=\frac52,
 \quad \Delta S=-8.
 \label{vii:vii:eq:valores-ejemplo-no-afin}
\end{equation}
L'expression affine \(-k_*s_*/4=-2\) diffère de la correction
exacte. Sur la branche positive,
\begin{equation}
 \partial_\lambda\Delta S=-\frac45,
 \qquad \partial_\lambda k_* =\frac1{20}
 \quad\text{dans }\lambda=25/2.
 \label{vii:vii:eq:derivadas-ejemplo-no-afin}
\end{equation}
\end{proposition}

\begin{proof}
Le gramien est \(1+k^2>0\), de sorte que le minimum possède les expressions indiquées. Différentier \(F\) donne \(s=-2\partial_kF\), puis \(\mathcal R=k-s/2\). Pour \(\lambda=25/2\), l'équation est \(k[(1+k^2)^2-25]=0\) ; puisque \(1+k^2>0\), ses racines sont précisément \(0,\pm2\). La hessienne totale est
\[
 1-\frac{2\lambda(1-3k^2)}{(1+k^2)^3},
\]
qui vaut \(16/5\) aux deux racines non nulles. La correction est
\(2+5/2-25/2=-8\). Dans l'identité radiale,
\[
 \int_0^1 k_*s(\lambda,tk_*)\,dt
    =\int_0^1\frac{200t}{(1+4t^2)^2}\,dt=20,
\]
car une primitive est \(-25/(1+4t^2)\).
Ainsi, \(\Delta S=8/4-20/2=-8\), y compris le terme non affine.
Pour \(\lambda>1/2\), la branche positive satisfait
\(k_*^2=\sqrt{2\lambda}-1\) et
\(\Delta S=\sqrt{2\lambda}-1/2-\lambda\).
Leurs dérivées donnent \eqref{vii:vii:eq:derivadas-ejemplo-no-afin}.
La formule d'enveloppe fournit le même résultat :
\((\partial_\lambda F)(\lambda,k_*)
 -(\partial_\lambda F)(\lambda,0)=1/5-1=-4/5\).
\end{proof}
